\section{Introducción}

El presente informe corresponde a la segunda etapa del Proyecto Civil, y tiene 
como propósito desarrollar el prediseño estructural del edificio residencial 
ubicado en Antofagasta, Chile. El proyecto consiste en un edificio de uso 
habitacional de 21 niveles sobre el terreno más un subterráneo destinado a 
estacionamientos, con una altura total de 53,42~m y un sistema estructural 
basado en muros de corte de hormigón armado, emplazado en Zona Sísmica~3 
sobre un suelo clasificado preliminarmente como Tipo~A.

El prediseño estructural permite definir las dimensiones preliminares de los 
elementos resistentes —losas, vigas y muros— de forma coherente con los 
criterios normativos, las demandas sísmicas y las restricciones arquitectónicas 
del proyecto. A diferencia del diseño definitivo, el prediseño se apoya en 
criterios simplificados y métodos de estimación directa que, si bien no 
reemplazan el análisis estructural riguroso, entregan una primera aproximación 
confiable de las dimensiones requeridas. Esta etapa resulta especialmente 
relevante en estructuras de mediana y gran altura como la presente, donde las 
decisiones tempranas de estructuración tienen un impacto directo en el 
comportamiento sísmico del edificio y en la eficiencia del uso de materiales.

La condición de alta sismicidad del emplazamiento —Zona Sísmica~3 según la 
NCh433— impone restricciones importantes sobre el diseño de los muros 
estructurales, los cuales deben proveer suficiente rigidez y resistencia lateral 
para controlar los desplazamientos y absorber las fuerzas sísmicas de manera 
adecuada. Al mismo tiempo, las restricciones arquitectónicas condicionan la 
disposición y dimensiones de vigas y losas, lo que hace necesario compatibilizar 
el prediseño estructural con los planos de arquitectura desde etapas tempranas.

Para ello, se desarrollan tres estudios principales: la determinación del espesor 
de losas por piso a partir de las losas de mayor luz, el predimensionamiento de 
vigas compatibilizado con los planos de arquitectura, y la estimación de espesores 
mínimos de muros estructurales mediante el método de densidad de muros aplicado 
en la zona crítica del edificio. Los resultados de cada estudio se presentan junto 
con los criterios de prediseño adoptados y las conclusiones derivadas del proceso 
de estructuración.